\scanpage{55}
% archive sheet number: 54

\uline{Corollaire} Dans tous les cas, la distribution des signes $\varepsilon_a(F)$ détermine celle des $\varepsilon_{a'}(F')$. \deleted{C'est clair \unsure{d}}

C'est clair d'abord dans les cas 1° et 3°, où on sait laquelle des faces $F,F'$ est la spécialisée de l'autre, \deleted{et par suite} Mais dans le cas \inserted{``symétrique''} 2°), \deleted{comme on connaît les signes $\varepsilon_a(F)$, on} On voit que \deleted{les} deux \inserted{grandes} arêtes extrêmes correspondantes pour $F,F'$ ont $=$ signe (c'est \unsure{tjrs} le cas, y compris dans les cas 1° et 2°), et pour les arêtes qui \inserted{ne} sont pas de ces types, elles sont (pour $F/F'$ fixé) toutes de $=$ signe (\uline{NB} leur nb, suivant les cas, est $(1,0)$, $(1,1)$ ou $(1,2)$) et les signes en question pour $F$ et pour $F'$ sont opposés.

\uline{Corollaire} Quand on connaît les signes $\varepsilon_F(a)$ pour \uline{une} face $F$, on les connaît pour \uline{ttes} (grâce à la connexité de $\mathcal X$).

J'ai supposé implicit\unsure{nt} \deleted{que} dans la vérification de la proposition que ni $F$, ni $F'$ n'est une face exceptionnelle (correspondant à une position relative exceptionnelle, i.e. ``supercritique''). Il faudrait que j'examine séparément \deleted{\unsure{de}} cas de telles faces. Pour une telle face, il n'y a que quatre grandes arêtes, et la distribution des signes est connue ---

Cela semble donner une façon simple de déduire les signes \inserted{(i.e. les \unsure{visées})} de \deleted{la} seule structure de carte de $(\mathcal X,K,S_K)$, sans avoir à utiliser les arcs critiques, et encore moins \deleted{les} le \deleted{\unsure{complexe}} graphe $L\subset\mathcal X$.
